\chapter{Install Git and configure SSH}
\label{ch:git-ssh-setup}

\noindent\textbf{Choose one Git installation method.} Install Git from the official website or use the terminal commands below. If Git is already installed, skip both methods and verify it with \texttt{git --version}.

\subsection{Option 1: Install Git from the official website}

\begin{tcolorbox}[
  enhanced,
  colback=white,
  colframe=UTPBlue!28,
  colbacktitle=UTPNavy,
  boxrule=0.65pt,
  arc=2mm,
  title={Choose your operating system},
  fonttitle=\bfseries\color{white},
  left=3.5mm,right=3.5mm,top=2.5mm,bottom=2.5mm
]
{\small\color{UTPMuted}Follow Git's official installation instructions for your platform.}\par\smallskip
\noindent
\begin{minipage}[t]{0.32\linewidth}
  \begin{tcolorbox}[enhanced,width=\linewidth,height=2.8cm,valign=center,colback=UTPPaleBlue!35,colframe=UTPBlue!30,boxrule=0.55pt,arc=1.5mm,left=1mm,right=1mm,top=1.5mm,bottom=1.5mm]
    \centering
    {\fontsize{23}{25}\selectfont\color{UTPBlue}\faWindows}\par\smallskip
    {\large\bfseries\color{UTPNavy}Windows}\par\smallskip
    \href{https://git-scm.com/install/windows}{\textbf{\small\color{UTPBlue}Installation guide $\rightarrow$}}
  \end{tcolorbox}
\end{minipage}\hfill
\begin{minipage}[t]{0.32\linewidth}
  \begin{tcolorbox}[enhanced,width=\linewidth,height=2.8cm,valign=center,colback=UTPPaleBlue!35,colframe=UTPBlue!30,boxrule=0.55pt,arc=1.5mm,left=1mm,right=1mm,top=1.5mm,bottom=1.5mm]
    \centering
    {\fontsize{23}{25}\selectfont\color{UTPBlue}\faApple}\par\smallskip
    {\large\bfseries\color{UTPNavy}macOS}\par\smallskip
    \href{https://git-scm.com/install/mac}{\textbf{\small\color{UTPBlue}Installation guide $\rightarrow$}}
  \end{tcolorbox}
\end{minipage}\hfill
\begin{minipage}[t]{0.32\linewidth}
  \begin{tcolorbox}[enhanced,width=\linewidth,height=2.8cm,valign=center,colback=UTPPaleBlue!35,colframe=UTPBlue!30,boxrule=0.55pt,arc=1.5mm,left=1mm,right=1mm,top=1.5mm,bottom=1.5mm]
    \centering
    {\fontsize{23}{25}\selectfont\color{UTPBlue}\faLinux}\par\smallskip
    {\large\bfseries\color{UTPNavy}Linux}\par\smallskip
    \href{https://git-scm.com/install/linux}{\textbf{\small\color{UTPBlue}Installation guide $\rightarrow$}}
  \end{tcolorbox}
\end{minipage}
\end{tcolorbox}

\subsection{Option 2: Install Git from the terminal}

Use this option instead of the website instructions above. Run the command for your operating system, then reopen the terminal:

\begin{lstlisting}[style=utpterminal]
# Windows PowerShell
winget install --id Git.Git -e --source winget

# macOS (Option A): install using Homebrew
brew install git

# macOS (Option B - If you don't have brew installed): install Apple's command-line tools
xcode-select --install

# Linux
sudo apt update
sudo apt install git
\end{lstlisting}

For SSH setup, use \textbf{PowerShell} on Windows or \textbf{Terminal} on macOS and Linux. The commands are labeled by platform. After either installation method, confirm Git is available:

\begin{lstlisting}[style=utpterminal]
git --version
\end{lstlisting}

\subsection{Configure your commit identity}

Set the name and email that should appear on future commits. Use an email
connected to your GitHub account, or your GitHub-provided private email.

\begin{lstlisting}[style=utpterminal]
git config --global user.name "Your Name"
git config --global user.email "you@example.com"
git config --global init.defaultBranch main

git config --global --list
\end{lstlisting}

\noindent\textbf{Choose one SSH setup option.} Option 1 uses video walkthroughs; Option 2 provides written terminal instructions. Both cover the SSH setup, so follow the path you prefer and skip steps already completed.

\subsection{Option 1: SSH setup with video walkthroughs}

\begin{tcolorbox}[
  enhanced,
  colback=white,
  colframe=UTPBlue!28,
  colbacktitle=UTPNavy,
  boxrule=0.65pt,
  arc=2mm,
  title={Choose your operating system},
  fonttitle=\bfseries\color{white},
  left=3.5mm,right=3.5mm,top=2.5mm,bottom=2.5mm
]
{\small\color{UTPMuted}Open the matching video guide for SSH setup.}\par\smallskip
\noindent
\begin{minipage}[t]{0.32\linewidth}
  \begin{tcolorbox}[enhanced,width=\linewidth,height=2.8cm,valign=center,colback=UTPPaleBlue!35,colframe=UTPBlue!30,boxrule=0.55pt,arc=1.5mm,left=1mm,right=1mm,top=1.5mm,bottom=1.5mm]
    \centering
    {\fontsize{23}{25}\selectfont\color{UTPBlue}\faWindows}\par\smallskip
    {\large\bfseries\color{UTPNavy}Windows}\par\smallskip
    \href{https://youtu.be/YT5YimhyX18?t=5}{\textbf{\small\color{UTPBlue}View video guide $\rightarrow$}}
  \end{tcolorbox}
\end{minipage}\hfill
\begin{minipage}[t]{0.32\linewidth}
  \begin{tcolorbox}[enhanced,width=\linewidth,height=2.8cm,valign=center,colback=UTPPaleBlue!35,colframe=UTPBlue!30,boxrule=0.55pt,arc=1.5mm,left=1mm,right=1mm,top=1.5mm,bottom=1.5mm]
    \centering
    {\fontsize{23}{25}\selectfont\color{UTPBlue}\faApple}\par\smallskip
    {\large\bfseries\color{UTPNavy}macOS}\par\smallskip
    \href{https://youtu.be/ACvE5Uy7ZB0?t=437}{\textbf{\small\color{UTPBlue}View video guide $\rightarrow$}}
  \end{tcolorbox}
\end{minipage}\hfill
\begin{minipage}[t]{0.32\linewidth}
  \begin{tcolorbox}[enhanced,width=\linewidth,height=2.8cm,valign=center,colback=UTPPaleBlue!35,colframe=UTPBlue!30,boxrule=0.55pt,arc=1.5mm,left=1mm,right=1mm,top=1.5mm,bottom=1.5mm]
    \centering
    {\fontsize{23}{25}\selectfont\color{UTPBlue}\faLinux}\par\smallskip
    {\large\bfseries\color{UTPNavy}Linux}\par\smallskip
    \href{https://youtu.be/AOM8CEyu3Ug?t=121}{\textbf{\small\color{UTPBlue}View video guide $\rightarrow$}}
  \end{tcolorbox}
\end{minipage}
\end{tcolorbox}

\subsection{Option 2: SSH setup from the terminal}

SSH lets Git authenticate without asking for a GitHub password on every push.
The key-generation process is similar across operating systems. Use PowerShell on Windows or Terminal on macOS and Linux. The \texttt{ssh-keygen} command is the same; the commands for checking existing keys and loading them into the SSH agent are shown separately by platform.

First check whether this computer already has a key:

\begin{lstlisting}[style=utpterminal]
# Windows PowerShell
Get-ChildItem -Force "$env:USERPROFILE\.ssh"

# macOS/Linux Terminal
ls -al ~/.ssh
\end{lstlisting}

If \texttt{id\_ed25519} and \texttt{id\_ed25519.pub} do not exist, create them.
Enter the email connected to GitHub and choose a passphrase when prompted:

\begin{lstlisting}[style=utpterminal]
ssh-keygen -t ed25519 -C "you@example.com"
\end{lstlisting}

The command above displays output similar to this. The save path and randomart may differ on your system:

\begin{enumerate}[leftmargin=1.7em]
  \item Generating public/private ed25519 key pair.
  \item Enter file in which to save the key (\textasciitilde/.ssh/id\_ed25519): \textit{(Press Enter to accept the default)}
  \item Enter passphrase (empty for no passphrase): \textit{(Press Enter for no passphrase)}
  \item Enter same passphrase again: \textit{(Press Enter to confirm)}
\end{enumerate}

\begin{lstlisting}[style=utpterminal]
Your identification has been saved in C:\Users\USERNAME\.ssh\id_ed25519
Your public key has been saved in C:\Users\USERNAME\.ssh\id_ed25519.pub
The key's randomart fingerprint is:
+--[ED25519 256]--+
|   ..oo+.        |
|   o.+.= .       |
|  . o o =        |
|   . E o         |
|    o . S        |
|     o o         |
|      o          |
+-----------------+
\end{lstlisting}

\pagebreak
\noindent\textbf{Load the private key into the SSH agent.} Choose the command for your operating system.

For Windows PowerShell:

\begin{lstlisting}[style=utpterminal]
Set-Service -Name ssh-agent -StartupType Automatic
Start-Service ssh-agent
ssh-add $env:USERPROFILE\.ssh\id_ed25519
git config --global core.sshCommand "C:/Windows/System32/OpenSSH/ssh.exe"
\end{lstlisting}

For macOS:
\begin{lstlisting}[style=utpterminal]
eval "$(ssh-agent -s)"
ssh-add --apple-use-keychain ~/.ssh/id_ed25519
\end{lstlisting}

For Linux:
\begin{lstlisting}[style=utpterminal]
eval "$(ssh-agent -s)"
ssh-add ~/.ssh/id_ed25519
\end{lstlisting}

Copy the \emph{public} key. Never upload or share the file without
\texttt{.pub}.

\begin{lstlisting}[style=utpterminal]
# Windows PowerShell
Get-Content $env:USERPROFILE\.ssh\id_ed25519.pub | Set-Clipboard

# macOS
pbcopy < ~/.ssh/id_ed25519.pub

# Linux: display it, then copy the complete line
cat ~/.ssh/id_ed25519.pub
\end{lstlisting}

\noindent To authenticate Git operations securely without entering your password, copy your public SSH key and register it in your GitHub account:

\begin{enumerate}[leftmargin=1.7em]
  \item Copy your public key directly to your clipboard. If you generated an ed25519 key, run the command below in PowerShell or terminal. Do not share or copy the private key (\texttt{id\_ed25519}).
  \item In the upper-right corner of GitHub, click your profile icon and select \textbf{Settings}.
  \item In the left sidebar under \textbf{Access}, select \textbf{SSH and GPG keys}.
  \item Click \textbf{New SSH key} (or \textbf{Add SSH key}).
  \item In the \textbf{Title} field, enter a descriptive label such as \texttt{personal laptop}.
  \item Ensure \textbf{Key type} is set to \textbf{Authentication Key}.
  \item Paste the public key string into the \textbf{Key} field.
  \item Click \textbf{Add SSH key}, then confirm your GitHub account password or two-factor prompt if requested.
\end{enumerate}

\noindent Verify that your machine successfully connects and authenticates with GitHub:

\begin{lstlisting}[style=utpterminal]
ssh -T git@github.com
\end{lstlisting}
